\documentclass[11pt]{article}
\usepackage{amsmath,amssymb,booktabs,graphicx}
\usepackage[margin=1in]{geometry}

\title{Sprint 1 Bootstrap: 36-State Factored MDP with LLM-Estimated Transitions}
\author{William Dennis}
\date{July 2026}

\begin{document}
\maketitle

\section{MDP Formulation}

The Sprint 1 simulator implements a finite-horizon factored MDP $(\mathcal{S}, \mathcal{A}, P, r, \gamma, T)$ with 36 states, 4 actions, and transitions estimated by an LLM (deepseek-v4-flash) using Algorithm 2.

\subsection{State Space}

Four factors: step~bin ($\texttt{step\_bin} \in \{\texttt{inactive},\ \texttt{moderate},\ \texttt{active}\}$), sleep quality ($\texttt{sleep} \in \{\texttt{good},\ \texttt{poor}\}$), day type ($\texttt{day\_of\_week} \in \{\texttt{weekday},\ \texttt{weekend}\}$), and burden level ($\texttt{burden} \in \{\texttt{low},\ \texttt{medium},\ \texttt{high}\}$):

\[
\mathcal{S} = \{\texttt{inactive},\ \texttt{moderate},\ \texttt{active}\}
\times \{\texttt{good},\ \texttt{poor}\}
\times \{\texttt{weekday},\ \texttt{weekend}\}
\times \{\texttt{low},\ \texttt{medium},\ \texttt{high}\}
\]

Three factors are deterministic: $\texttt{day\_of\_week}$ cycles daily through a 7-day pattern; $\texttt{burden}$ uses a rolling window of the last 3 actions (increments on non-idle actions). Two factors are stochastic: $\texttt{step\_bin}$ at each timestep, and $\texttt{sleep}$ at day boundaries.

\subsection{Action Space}

\[
\mathcal{A} = \{\texttt{idle},\ \texttt{movement\_suggestion},\ \texttt{goal\_reminder},\ \texttt{journal}\}
\]

All non-idle actions incur a penalty of 0.05.

\subsection{Transition Function}

Transitions are loaded from 6 JSON probability tables generated by an LLM (deepseek-v4-flash, $N = 16$ samples per state-action pair, Algorithm 2):

\begin{itemize}
  \item \texttt{day\_boundary.json}: $P(\texttt{sleep}' \mid \texttt{step\_bin},\ \texttt{burden},\ \texttt{day\_of\_week},\ \texttt{sleep})$ -- action-independent, evaluated at step~0 of each day.
  \item \texttt{within\_day\_0.json} through \texttt{within\_day\_4.json}: $P(\texttt{step\_bin}' \mid \texttt{step\_bin},\ \texttt{burden},\ \texttt{action},\ \texttt{day\_of\_week},\ \texttt{sleep})$ -- one table per timestep, action-conditioned.
\end{itemize}

The LLM's estimates reflect plausible real-world dynamics: interventions have modest effects on step counts, and sleep quality influences next-day activity.

\subsection{Reward}

\[
R = \alpha \cdot \texttt{step\_bin\_value} + (1 - \alpha) \cdot \texttt{sleep\_value} - \texttt{action\_penalty}
\]

with $\alpha = 0.9$. Step~bin values: $\texttt{inactive} \to 0.0$, $\texttt{moderate} \to 0.5$, $\texttt{active} \to 1.0$. Sleep values: $\texttt{good} \to 1.0$, $\texttt{poor} \to -1.0$. Action penalty: 0.0 for idle, 0.05 otherwise.

Optional \texttt{reward\_multiplier\_by\_step} masks the last step of each day: $[1, 1, 1, 1, 0]$.

\subsection{Episode}

$T = 90$ days $\times$ 5 steps/day $= 450$ timesteps per episode.

\section{Reference Policy Bounds}

Six reference policies evaluated via simulation (10 seeds, 450 timesteps):

\begin{table}[h]
\centering
\begin{tabular}{lrrr}
\toprule
Policy & Total Reward & Per Step & Description \\
\midrule
Fixed idle & 380.4 $\pm$ 3.7 & 0.845 & Never intervene (no penalty) \\
Greedy oracle (1-step) & 281.0 $\pm$ 20.9 & 0.625 & Myopic expected-reward maximisation \\
Random & 100.1 $\pm$ 11.7 & 0.222 & Uniform action selection \\
Fixed movement\_suggestion & 35.2 $\pm$ 4.8 & 0.078 & Always suggest movement \\
Fixed journal & 50.5 $\pm$ 5.1 & 0.112 & Always journal \\
Fixed goal\_reminder & 91.2 $\pm$ 7.1 & 0.203 & Always remind of goal \\
\bottomrule
\end{tabular}
\caption{Reference policy bounds (10 seeds). Always-idle dominates because the LLM estimates natural activity is already high and action penalties are costly. The greedy oracle's myopic one-step lookahead fails to capture long-term value.}
\end{table}

Always-idle at 0.845/step sets an upper bound for agents: the LLM estimates that most users maintain moderate-to-active step counts without intervention, and paying a 0.05 penalty for no improvement is worse than doing nothing.

\section{Base Benchmark Results}

50 seeds, 450 timesteps, default hyperparameters (EG $\epsilon=0.1$, UCB $c=2.0$, TS $\alpha_0=\beta_0=1$). Config: \texttt{configs/sprint1\_bootstrap.yaml}.

\begin{table}[h]
\centering
\begin{tabular}{lrrr}
\toprule
Agent & Total Reward & Per Step & Last 50 Steps \\
\midrule
Epsilon-Greedy ($\epsilon=0.1$) & 280.8 $\pm$ 89.3 & 0.624 & 0.743 \\
Thompson Sampling ($\alpha_0=\beta_0=1$) & 252.6 $\pm$ 105.6 & 0.561 & 0.662 \\
UCB ($c=2.0$) & 240.4 $\pm$ 33.9 & 0.534 & 0.803 \\
Random & 99.7 $\pm$ 10.9 & 0.222 & 0.215 \\
\bottomrule
\end{tabular}
\caption{Base benchmark results (50 seeds, 450 timesteps). Learning agents massively outperform random (2.5--2.8$\times$), confirming the bootstrap tables contain exploitable structure. EG leads.}
\end{table}

Key observation: under bootstrap transitions, learning agents achieve 2.5--2.8$\times$ random, starkly different from the random-transition baseline where all agents clustered within 3\% of each other. The LLM-estimated dynamics have meaningful structure.

\section{Full Benchmark Results}

50 seeds, 450 timesteps, 9 agents with MVP-optimised hyperparameters (EG $\epsilon=0.05$, D-EG $\epsilon_{\text{start}}=0.2$, $T_{\text{decay}}=200$, UCB $c=0.5$). Config: \texttt{configs/sprint1\_bootstrap\_extensions.yaml}.

\begin{table}[h]
\centering
\begin{tabular}{lrrr}
\toprule
Agent & Total Reward & Per Step & Last 50 Steps \\
\midrule
Standard UCB ($c=0.5$) & 358.1 $\pm$ 14.0 & 0.796 & 0.837 \\
Standard D-EG & 287.8 $\pm$ 106.9 & 0.640 & 0.718 \\
Standard EG ($\epsilon=0.05$) & 282.3 $\pm$ 95.0 & 0.627 & 0.748 \\
Contextual UCB & 268.3 $\pm$ 29.6 & 0.596 & 0.733 \\
Standard TS & 252.6 $\pm$ 105.6 & 0.561 & 0.662 \\
Contextual D-EG & 240.4 $\pm$ 93.0 & 0.534 & 0.614 \\
Contextual EG & 216.9 $\pm$ 95.7 & 0.482 & 0.583 \\
Contextual TS & 155.5 $\pm$ 70.8 & 0.346 & 0.399 \\
Random & 99.7 $\pm$ 10.9 & 0.222 & 0.215 \\
\midrule
Fixed idle (upper bound) & 380.4 & 0.845 & 0.845 \\
\bottomrule
\end{tabular}
\caption{Full benchmark results (50 seeds). Standard UCB at $c=0.5$ reaches 94\% of the always-idle ceiling. Contextual agents universally underperform their standard counterparts.}
\end{table}

Key findings:
\begin{enumerate}
  \item \textbf{Std UCB at $c=0.5$ is dominant}: 358.1 total, 94\% of the always-idle ceiling, with the lowest variance (14.0 vs 70--106 for others). Hyperparameter tuning is critical: the default $c=2.0$ yielded only 240.4.
  \item \textbf{Contextual agents are worse}: Every contextual variant underperforms its standard counterpart, the opposite of the MVP result. This suggests $\texttt{step\_bin}$ as a context feature misleads agents in the bootstrap dynamics -- transitions are non-stationary (vary by timestep) and the current $\texttt{step\_bin}$ does not reliably predict intervention efficacy.
  \item \textbf{The $c=0.5$ UCB achieves consistency}: with $\pm$14.0 std dev vs $\pm$33.9 at $c=2.0$, tuning simultaneously improves both mean and variance.
\end{enumerate}

\section{Masked Reward Benchmark}

Same 9 agents on \texttt{configs/sprint1\_bootstrap\_extensions\_masked.yaml} where step~4 (end of day) yields zero reward. Results are 50 seeds.

\begin{table}[h]
\centering
\begin{tabular}{lrrr}
\toprule
Agent & Total Reward & Per Step & Last 50 Steps \\
\midrule
Standard UCB ($c=0.5$) & 265.9 $\pm$ 37.3 & 0.591 & 0.681 \\
Standard D-EG & 248.3 $\pm$ 76.2 & 0.552 & 0.616 \\
Standard EG ($\epsilon=0.05$) & 195.8 $\pm$ 98.4 & 0.435 & 0.518 \\
Contextual UCB & 178.4 $\pm$ 35.2 & 0.397 & 0.501 \\
Standard TS & 174.6 $\pm$ 70.6 & 0.388 & 0.456 \\
Contextual D-EG & 169.9 $\pm$ 70.0 & 0.378 & 0.435 \\
Contextual EG & 163.7 $\pm$ 73.9 & 0.364 & 0.443 \\
Contextual TS & 119.0 $\pm$ 38.5 & 0.264 & 0.289 \\
Random & 83.6 $\pm$ 8.4 & 0.186 & 0.179 \\
\bottomrule
\end{tabular}
\caption{Masked benchmark results (50 seeds, step~4 zero-reward). Relative ordering is preserved. D-EG is most mask-resilient (86\% retention).}
\end{table}

\section{Comparison to Random-Transition Baseline}

\begin{table}[h]
\centering
\begin{tabular}{lrr}
\toprule
Metric & Random Transitions & Bootstrap Transitions \\
\midrule
Best agent & Std D-EG 189.1 (0.420) & Std UCB 358.1 (0.796) \\
Random baseline & 183.5 (0.408) & 99.7 (0.222) \\
Best vs random gap & 3\% & 259\% \\
Contextual advantage & None ($\approx$0\%) & Negative (all worse) \\
Variance (best agent) & $\pm$82 & $\pm$14 \\
\bottomrule
\end{tabular}
\caption{Bootstrap transitions vs random transitions. The LLM-estimated tables produce structured dynamics with large agent differentials, while random transitions wash out all policy differences. Std UCB's low variance under bootstrap suggests reliable convergence.}
\end{table}

\section{Context Feature Comparison}

The extensions benchmark used $\texttt{step\_bin}$ as the context feature (3 bins, same fragmentation as step\_bin). To test whether a better-chosen context could close the gap, we swapped to $\texttt{burden}$ (3 levels, same fragmentation) and re-ran the 9-agent benchmark. Results are 50 seeds.

\begin{table}[h]
\centering
\begin{tabular}{lrrr}
\toprule
Agent & No Context & Ctx $\texttt{step\_bin}$ & Ctx $\texttt{burden}$ \\
\midrule
Standard UCB ($c=0.5$) & 358.1 $\pm$ 14.0 & --- & --- \\
Contextual UCB & --- & 268.3 $\pm$ 29.6 & 250.6 $\pm$ 39.7 \\
Contextual D-EG & --- & 240.4 $\pm$ 93.0 & 216.0 $\pm$ 77.6 \\
Contextual EG & --- & 216.9 $\pm$ 95.7 & 200.4 $\pm$ 82.9 \\
Contextual TS & --- & 155.5 $\pm$ 70.8 & 158.3 $\pm$ 56.1 \\
\midrule
Random & 99.7 $\pm$ 10.9 & --- & --- \\
\bottomrule
\end{tabular}
\caption{Comparison of context features (50 seeds). Burden context underperforms step\_bin context across all agents except TS (where the difference is negligible).}
\end{table}

Despite having 13$\times$ higher ANOVA F-statistic as a predictor of next-step activity, burden makes a worse context feature. The reason is that $\texttt{burden}$ is \emph{endogenous} -- it is a rolling window count of the agent's past non-idle actions. When a fragmented contextual agent explores, it takes non-idle actions, which increase burden, which shifts the context distribution, which triggers further exploration. This creates a feedback loop:

\[
\text{exploration} \rightarrow \text{higher burden} \rightarrow \text{unfamiliar context} \rightarrow \text{more exploration} \rightarrow \cdots
\]

In contrast, $\texttt{step\_bin}$ is at least partly exogenous -- it depends on the LLM-estimated transition dynamics and is not directly controlled by recent actions. The step\_bin distribution is relatively stable regardless of the agent's policy, making it a safer (if less predictive) context feature.

Both context features are dominated by the standard (non-contextual) UCB, confirming that $\texttt{step\_bin}$ context does not help this MDP.

\section{Horizon Scaling}

To test whether contextual agents would catch the standard UCB given more data, we re-ran both context-feature benchmarks at 900 days (4500 steps, 10$\times$) and 9000 days (45000 steps, 100$\times$). All other settings are identical (50 seeds, same hyperparameters).

\begin{table}[h]
\centering
\begin{tabular}{lrrrrr}
\toprule
Agent & 90d / 450 st. & 900d / 4500 st. & 9000d / 45000 st. & Per-step \\
\midrule
Std UCB ($c=0.5$) & 358.1 $\pm$ 14.0 & 3795.7 $\pm$ 24.0 & 38241.5 $\pm$ 51.1 & 0.8498 \\
Ctx UCB ($\texttt{step\_bin}$) & 268.3 $\pm$ 29.6 & 3606.6 $\pm$ 43.5 & 37971.9 $\pm$ 63.9 & 0.8438 \\
Ctx UCB ($\texttt{burden}$) & 250.6 $\pm$ 39.7 & 3389.9 $\pm$ 79.5 & 36404.1 $\pm$ 374.9 & 0.8090 \\
Std EG ($\epsilon=0.05$) & 282.3 $\pm$ 95.0 & 3578.6 $\pm$ 264.7 & 37182.9 $\pm$ 264.0 & 0.8263 \\
Ctx EG ($\texttt{step\_bin}$) & 216.9 $\pm$ 95.7 & 3376.6 $\pm$ 423.0 & 36824.1 $\pm$ 952.6 & 0.8183 \\
Std D-EG & 287.8 $\pm$ 106.9 & 3373.3 $\pm$ 916.3 & 37097.9 $\pm$ 2987.5 & 0.8244 \\
Ctx D-EG ($\texttt{step\_bin}$) & 240.4 $\pm$ 93.0 & 3168.0 $\pm$ 763.3 & 37057.1 $\pm$ 1590.5 & 0.8235 \\
Std TS & 252.6 $\pm$ 105.6 & 3354.3 $\pm$ 842.0 & 37060.1 $\pm$ 3511.0 & 0.8236 \\
Ctx EG ($\texttt{burden}$) & 200.4 $\pm$ 82.9 & 3140.1 $\pm$ 386.4 & 35008.6 $\pm$ 1722.7 & 0.7780 \\
Ctx D-EG ($\texttt{burden}$) & 216.0 $\pm$ 77.6 & 2694.5 $\pm$ 927.8 & 31957.4 $\pm$ 6113.1 & 0.7102 \\
Ctx TS ($\texttt{step\_bin}$) & 155.5 $\pm$ 70.8 & 2215.5 $\pm$ 737.6 & 25537.2 $\pm$ 8989.0 & 0.5675 \\
Ctx TS ($\texttt{burden}$) & 158.3 $\pm$ 56.1 & 1776.2 $\pm$ 571.3 & 19794.5 $\pm$ 5612.9 & 0.4399 \\
\midrule
Random & 99.7 $\pm$ 10.9 & 1020.8 $\pm$ 45.1 & 10177.4 $\pm$ 118.4 & 0.2262 \\
\bottomrule
\end{tabular}
\caption{Horizon scaling: 90d, 900d, and 9000d (50 seeds). At 9000d Ctx UCB (step\_bin) reaches 99.3\% of Std UCB. Rank ordering is stable across all three horizons.}
\end{table}

Several patterns emerge:

\begin{enumerate}
  \item \textbf{Ctx UCB ($\texttt{step\_bin}$) effectively closes the gap}. At 90d / 450 steps it achieves 75\% of Std UCB's total; at 900d / 4500 steps it reaches 95\%; at 9000d / 45000 steps it reaches 99.3\% (37971.9 vs 38241.5). The contextual fragmentation penalty is purely a \emph{sample-efficiency cost}, not an asymptotic limitation. Given enough data, step\_bin-contextual UCB converges to the same optimal policy of doing nothing.

  \item \textbf{Ctx UCB ($\texttt{burden}$) also converges, but slower}. Even at 45000 steps, burden-context UCB reaches only 0.8090/step vs Std UCB's 0.8498 (95.2\%). The endogenous feedback loop does not disappear with more data, but its relative penalty shrinks: 75\% $\to$ 89\% $\to$ 95\% as horizon increases.

  \item \textbf{The 10$\times$ multiplier gap disappears at longer horizons}. At 900d, contextual agents had higher gain multipliers (13--16$\times$) than standard ones (11--13$\times$). At 9000d the per-step rates show a different pattern: Std UCB converges to 0.8498, with most standard agents clustering at 0.823--0.826 and step\_bin-contextual agents at 0.818--0.844. The gap is now driven by \emph{asymptotic convergence rate} rather than a fixed fragmentation penalty.

  \item \textbf{Rank ordering is remarkably stable across all three horizons}. No agent that was worse at 450 steps overtakes one that was better at 45000 steps. The relative hierarchy is fully determined early in training. However, some pairs converge: Ctx D-EG ($\texttt{step\_bin}$) nearly ties Std TS and Std D-EG at 45000 steps, suggesting the exploration decay eventually lets it converge.

  \item \textbf{Std UCB at 45000 steps reaches 0.8498/step, matching the always-idle upper bound (0.845)}. The standard deviation shrinks to $\pm$51.1, the lowest of any agent at any horizon. UCB with $c=0.5$ has converged to the optimal deterministic policy of always choosing idle.

  \item \textbf{Contextual TS remains terrible even at 45000 steps}. Ctx TS ($\texttt{step\_bin}$) at 0.5675/step and Ctx TS ($\texttt{burden}$) at 0.4399/step are the only agents that do not converge toward the standard-agent cluster. Beta-Bernoulli Thompson Sampling with independent per-context Beta posteriors cannot aggregate information across contexts, so it never escapes the fragmentation penalty.

  \item \textbf{$\texttt{step\_bin}$ beats $\texttt{burden}$ at all three horizons}. The gap between step\_bin and burden UCB narrows from 17.7 (90d) to 216.7 (900d) to 1567.8 (9000d), but in relative terms burden per-step stays at 95--96\% of step\_bin per-step at all three horizons. The relative penalty is constant.
\end{enumerate}

\section{Key Findings}

\begin{enumerate}
  \item \textbf{Bootstrap transitions have exploitable structure}. Unlike the Dirichlet-random baselines where all agents performed similarly ($\sim$185 total), the LLM-estimated tables yield a 259\% gap between best learner and random. The LLM produces non-trivial transition dynamics.
  \item \textbf{Always-idle is near-optimal}. The LLM estimates that most users naturally maintain moderate-to-active step counts. Interventions incur a penalty without consistently improving outcomes, making inaction the best policy. Std UCB learns this implicitly, reaching 94\% of the idle ceiling.
  \item \textbf{Context features hurt, even with better signal}. Swapping from $\texttt{step\_bin}$ to $\texttt{burden}$ -- a context with 13$\times$ stronger predictive power -- made performance worse, not better. The reason is that burden is endogenous: exploration changes burden, which shifts the context distribution, which triggers more exploration. This feedback loop amplifies the sample-efficiency cost of context fragmentation.
  \item \textbf{UCB hyperparameter sensitivity is extreme}. Default $c=2.0$: 240. Tuned $c=0.5$: 358. The MVP-optimised values transferred well, but a dedicated grid search may find even better settings.
  \item \textbf{The greedy oracle (one-step lookahead) is not a good bound} at 281, far below Std UCB's 358. Myopic expected-reward maximisation fails because actions affect future states through burden accumulation and sleep at day boundaries.
  \item \textbf{The LLM's transition estimates are a plausible starting point} but should be validated against real data. The strong idle preference may reflect LLM bias toward inaction rather than true user dynamics.
\end{enumerate}

\end{document}
